\section{Conclusion}

We set out to build a version control system that records the edits a developer
asks for instead of the lines an agent writes, and to test that system honestly.
The evaluation found that the read operations work and the write operations fail.
Attribution, preview, and dependency display survived first contact with a user
who had never seen the tool, and a participant in the baseline condition
independently described the capability the tool provides. The write operations
did not survive. On 27 of 28 repositories we did not record ourselves, the
primary verbs refused to run because they could not guarantee the rebuilt file
would match what the developer left on disk. The refusal is correct, but it
converts a tool that promises selective revert into a tool that declines to
record the next edit. The cause is the fork rule. The fork rule withdraws
records to protect 162 actually-forked functions. That withdrawal cascades
eighteen-fold to cover 2{,}864 function keys, the rebuild becomes incomplete, and
every write guard fires on the gap. The fix is straightforward in principle
(loosen the fork rule) but involves a real tradeoff: loosening it gives up the
safety guarantee of Section~\ref{sec:fork} in exchange for the completeness of
Section~\ref{sec:eval-exact}.

Section~\ref{sec:discussion-synthesis} traces each design commitment back to its
evaluation outcome. The pattern across all four commitments is the same: the part
that shows information to the developer held up under first contact, and the part
that modifies the developer's files either failed or refused to act.

Three findings generalise past this system. First, automating the earlier
stages of Parasuraman's model (parsing, clustering, dependency tracking)
delivers value independently of whether the later stages work, and a developer
who trusts the reports can use them with git directly. Second, the failure mode
that cost us the most was a command that reported success while doing nothing.
A silent lie is worse than a crash because a tool that lies quietly erodes the
basis for trusting anything else it says, and that erosion did not recover
within a session (Section~\ref{sec:pilot-lessons}). Third, evaluating intent
recovery against real agentic transcripts is unstable in a way the metric hides.
The share of contentless turns varies from 0\% to 88\% across our four
repositories, and a pairwise F1 computed on a history with 88\% continuation
tokens measures the ground truth's grain rather than the system's accuracy
(Section~\ref{sec:codebook-wrong}).

The second finding generalises further. Durability and legibility require
different kinds of guarantee. The append-only store ensures that every prior
state is reachable, and it does so structurally. But the developer still has to
know which command to type and whether its output told the truth, and the
architecture provides no structural guarantee of that. Every one of our seven
self-reporting failures demonstrated this gap: the data was present, and the
tool's account of how to reach it was wrong.
